% GENERATED FILE. Edit canonical lessons, then run npm run generate:books.
% canonical-source-hash: fnv1a64:f7d2e3f2c9ef0e0f
% canonical-lessons: ZH-C361-dazhaohu, ZH-C361-zhaodai, ZH-C361-jugong, ZH-C361-jingjiu, ZH-C361-kanbuqi

\chapter{A Good Guest, a Good Host}
\label{ch:zh-a-good-guest-a-good-host}
\hlchaptermodality{361}

\begin{chapteropening}
\textbf{By the end of this chapter:} \emph{I can say hello properly, bow, look after guests and propose a toast.}

Complete the last lesson of chapter 361: I can say hello properly, bow, look after guests and propose a toast.
\end{chapteropening}

\section[dǎ zhāohu]{dǎ zhāohu — to say hello to someone}

\label{lesson:ZH-C361-dazhaohu}



\begin{quote}
Before the new one: say the Mandarin for the occasion you are in, then the Mandarin for to jump the queue.
\end{quote}

\subsection*{You'll want to know: dǎ zhāohu}
\emph{dǎ zhāohu} — \textquotedblleft{}to say hello to someone\textquotedblright{}.

\begin{quote}
\emph{wǒmen jiànmiàn de shíhou yīnggāi dǎ zhāohu} — We should say hello when we meet.
\end{quote}

Its characters come later; for now it is a word for the ear.

\begin{cousinweb}
\emph{dǎ zhāohu} is built from \zh{打} 'to strike' and \emph{zhāohu}, all three inherited from Middle Chinese, \emph{taengX}, \emph{tsyew} and \emph{xu} (Guangyun, in Baxter's transcription), which glosses the first 'to strike', the second 'to beckon' and the third 'to call out'. The last syllable is said light.
\end{cousinweb}

\begin{grammarlens}[title={What you've built}]
One more word for giving reasons.
\end{grammarlens}

\subsection*{Your turn}
\begin{itemize}
\raggedright
  \item \emph{Say it:} \emph{dǎ zhāohu}
  \item \emph{Say it:} \emph{dǎ zhāohu}, once more
  \item \emph{Say it:} \emph{wǒmen jiànmiàn de shíhou yīnggāi dǎ zhāohu} — We should say hello when we meet.
  \item \emph{Recall:} say the Mandarin for the occasion you are in, then the Mandarin for to jump the queue, then say \emph{dǎ zhāohu} again
\end{itemize}

\begin{tcolorbox}[breakable,colback=teal!4,colframe=teal!35!black,title={Before you move on}]
What does dǎ zhāohu mean? (\textquotedblleft{}To say hello to someone\textquotedblright{}.) Say it once more.
\end{tcolorbox}
\section[zhāodài]{zhāodài — to look after your guests; to entertain}

\label{lesson:ZH-C361-zhaodai}



\begin{quote}
Before the new one: say the Mandarin for to jump the queue, then the Mandarin for to say hello to someone.
\end{quote}

\subsection*{You'll want to know: zhāodài}
\emph{zhāodài} — \textquotedblleft{}to look after your guests; to entertain\textquotedblright{}.

\begin{quote}
\emph{tāmen zhāodài kèrén hěn rèqíng} — They're very warm hosts.
\end{quote}

Its characters come later; for now it is a word for the ear.

\begin{cousinweb}
\emph{zhāodài} is built from \emph{zhāo} and \emph{dài}, both inherited from Middle Chinese, \emph{tsyew} and \emph{dojX} (Guangyun, in Baxter's transcription), which glosses the first 'to beckon' and the second 'to wait'.
\end{cousinweb}

\begin{grammarlens}[title={What you've built}]
One more word for giving reasons.
\end{grammarlens}

\subsection*{Your turn}
\begin{itemize}
\raggedright
  \item \emph{Say it:} \emph{zhāodài}
  \item \emph{Say it:} \emph{zhāodài}, once more
  \item \emph{Say it:} \emph{tāmen zhāodài kèrén hěn rèqíng} — They're very warm hosts.
  \item \emph{Recall:} say the Mandarin for to jump the queue, then the Mandarin for to say hello to someone, then say \emph{zhāodài} again
\end{itemize}

\begin{tcolorbox}[breakable,colback=teal!4,colframe=teal!35!black,title={Before you move on}]
What does zhāodài mean? (\textquotedblleft{}To look after your guests; to entertain\textquotedblright{}.) Say it once more.
\end{tcolorbox}
\section[jūgōng]{jūgōng — to bow}

\label{lesson:ZH-C361-jugong}



\begin{quote}
Before the new one: say the Mandarin for to say hello to someone, then the Mandarin for to look after your guests.
\end{quote}

\subsection*{You'll want to know: jūgōng}
\emph{jūgōng} — \textquotedblleft{}to bow\textquotedblright{}.

\begin{quote}
\emph{yǎnchū yǐhòu yǎnyuán yìqǐ jūgōng} — After the show, the actors bow together.
\end{quote}

Its characters come later; for now it is a word for the ear.

\begin{cousinweb}
\emph{jūgōng} is built from \emph{jū} and \emph{gōng}, both inherited from Middle Chinese, \emph{kjuwk} and \emph{kjuwng} (Guangyun, in Baxter's transcription), which glosses the first 'to examine closely' and 'to rear', and the second 'the body'. How the pair came to mean bowing is not settled.
\end{cousinweb}

\begin{grammarlens}[title={What you've built}]
One more word for giving reasons.
\end{grammarlens}

\subsection*{Your turn}
\begin{itemize}
\raggedright
  \item \emph{Say it:} \emph{jūgōng}
  \item \emph{Say it:} \emph{jūgōng}, once more
  \item \emph{Say it:} \emph{yǎnchū yǐhòu yǎnyuán yìqǐ jūgōng} — After the show, the actors bow together.
  \item \emph{Recall:} say the Mandarin for to say hello to someone, then the Mandarin for to look after your guests, then say \emph{jūgōng} again
\end{itemize}

\begin{tcolorbox}[breakable,colback=teal!4,colframe=teal!35!black,title={Before you move on}]
What does jūgōng mean? (\textquotedblleft{}To bow\textquotedblright{}.) Say it once more.
\end{tcolorbox}
\section[jìngjiǔ]{jìngjiǔ — to propose a toast}

\label{lesson:ZH-C361-jingjiu}



\begin{quote}
Before the new one: say the Mandarin for to look after your guests, then the Mandarin for to bow.
\end{quote}

\subsection*{You'll want to know: jìngjiǔ}
\emph{jìngjiǔ} — \textquotedblleft{}to propose a toast\textquotedblright{}.

\begin{quote}
\emph{zài Zhōngguó chī fàn de shíhou dàjiā chángcháng jìngjiǔ} — In China people often propose toasts at meals.
\end{quote}

Its characters come later; for now it is a word for the ear.

\begin{cousinweb}
\emph{jìngjiǔ} is built from \emph{jìng} and \emph{jiǔ}, both inherited from Middle Chinese, \emph{kjaengH} and \emph{tsjuwX} (Guangyun, in Baxter's transcription), which glosses the first 'respectful' and the second 'wine'.
\end{cousinweb}

\begin{grammarlens}[title={What you've built}]
One more word for giving reasons.
\end{grammarlens}

\subsection*{Your turn}
\begin{itemize}
\raggedright
  \item \emph{Say it:} \emph{jìngjiǔ}
  \item \emph{Say it:} \emph{jìngjiǔ}, once more
  \item \emph{Say it:} \emph{zài Zhōngguó chī fàn de shíhou dàjiā chángcháng jìngjiǔ} — In China people often propose toasts at meals.
  \item \emph{Recall:} say the Mandarin for to look after your guests, then the Mandarin for to bow, then say \emph{jìngjiǔ} again
\end{itemize}

\begin{tcolorbox}[breakable,colback=teal!4,colframe=teal!35!black,title={Before you move on}]
What does jìngjiǔ mean? (\textquotedblleft{}To propose a toast\textquotedblright{}.) Say it once more.
\end{tcolorbox}
\section[\texorpdfstring{\zh{看不起} (kànbuqǐ)}{看不起 (kànbuqǐ)}]{\zh{看不起} (kànbuqǐ) — to look down on}

\label{lesson:ZH-C361-kanbuqi}



\begin{quote}
Before the new one: say the Mandarin for to bow, then the Mandarin for to propose a toast.
\end{quote}

\subsection*{You'll want to know: \zh{看不起}}
\textbf{\zh{看不起}} — \emph{kànbuqǐ} — \textquotedblleft{}to look down on\textquotedblright{}.

\begin{quote}
\textbf{\zh{你不应该看不起他}} \emph{nǐ bù yīnggāi kànbuqǐ tā} — You shouldn't look down on him.
\end{quote}

The tone of bù is lost here: between the two parts it is said light, bu.

\begin{cousinweb}
\textbf{\zh{看不起}} (\emph{kànbuqǐ}) is built from \zh{看}, \zh{不} 'not' and \zh{起} 'to rise', all three inherited from Middle Chinese, \emph{khan}, \emph{pjut} and \emph{khiX} (Guangyun, in Baxter's transcription), which glosses the first 'to view' and the third 'to rise'. How the three came to mean looking down on someone is not settled.
\end{cousinweb}

\begin{grammarlens}[title={What you've built}]
One more word for giving reasons.
\end{grammarlens}

\subsection*{Your turn}
\begin{itemize}
\raggedright
  \item \emph{Say it:} \emph{kànbuqǐ}
  \item \emph{Say it:} \emph{kànbuqǐ}, once more
  \item \emph{Say it:} \emph{nǐ bù yīnggāi kànbuqǐ tā} — You shouldn't look down on him.
  \item \emph{Recall:} say the Mandarin for to bow, then the Mandarin for to propose a toast, then say \emph{kànbuqǐ} again
\end{itemize}

\begin{tcolorbox}[breakable,colback=teal!4,colframe=teal!35!black,title={Before you move on}]
What does \zh{看不起} mean? (\textquotedblleft{}To look down on\textquotedblright{}.) Say it once more.
\end{tcolorbox}
